\section{Заключение}
\label{sec:Chapter5} \index{Chapter5}

Обратимся к целям и задачам, сформулированным в начале работы, и проанализируем
полученные результаты.

Основной целью работы было расширение типовой системы и увеличение множества проверяемых
инструкций. Поставленная цель была
достигнута за счет грамотной поэтапной реализации решения.

В ходе работы были получены следующие результаты.
\begin{enumerate}
    \item Были изучены верификаторы байт-кода статически и динамически типизированных 
    виртуальных машин. Разобраны основные
    этапы проверки метода: построение
    графа потока управления, создание верификаторного контекста,
    абстрактная интерпретация инструкций. Рассмотрены свойства
    статической проверки:
    консервативность анализа, конечность алгоритма,
    согласование абстрактной типовой модели с семантикой времени исполнения.

    \item Проанализированы особенности plugin-расширения набора инструкций.
    Было уделено внимание инструкциям вызова метода по имени \newline
    \texttt{plugin.call.name.*}, которые используются в сценариях, где конкретная
    реализация метода зависит от фактического класса объекта во время
    исполнения.

    \item Проанализировано множество инструкций \texttt{any.*}. Было
    показано, что значение типа \texttt{any} может содержать как значение,
    пришедшее из динамической среды, так и значение из статической части
    программы. Поэтому проверка таких инструкций должна учитывать не только
    общий факт ссылочного представления, но и особенности операций доступа к
    свойствам, индексаторам, вызовам функций и проверкам типа на границе между
    статической и динамической семантикой.

    \item Была расширена типовая система верификатора. В неё был добавлен
    новый тип - union, представляющий собой объединение нескольких ссылочных типов.

    \item Реализованы верификаторные обработчики инструкций вызова метода по имени.
    При этом проверки, зависящие от состояния
    программы во время исполнения, такие как позднее разрешение метода по
    фактическому классу объекта, обработка отсутствия подходящей реализации и
    проверка абстрактного метода, были оставлены на стороне интерпретатора.

    \item Реализована верификация инструкций \texttt{any.*}. В работу
    вошли инструкции вызова функций и методов, вызова конструкторов, чтения и
    записи свойств по имени, индексу и значению-ключу, а также инструкция
    \texttt{any.isinstance}. Для этих инструкций также 
    были уточнены правила в ISA.

    \item Разработанные изменения были покрыты тестами: были добавлены 
    положительные и отрицательные сценарии,
    покрывающие разные формы передачи аргументов.
    Также была доработана тестовая инфраструктура для запуска
    тестов верификатора в управляемом контексте исполнения.
\end{enumerate}

Практическая значимость работы заключается в том, что верификатор стал
проверять больший класс программ: с плагиноспецифическими инструкциями и операциями над
значениями типа \texttt{any}. Это уменьшило число ситуаций, когда корректный
байт-код не может быть содержательно проверен до исполнения, и
повысило точность анализа за счет более полного соответствия между типовой
моделью верификатора и моделью времени исполнения платформы.

Таким образом, поставленные задачи были выполнены. Результаты работы
были добавлены в программный продукт.
